\documentclass[11pt]{article}
\usepackage[margin=1in]{geometry}
\usepackage[T1]{fontenc}
\usepackage{lmodern}
\usepackage{microtype}
\usepackage{booktabs}
\usepackage{enumitem}
\usepackage{hyperref}
\usepackage[numbers]{natbib}

\title{When Does Deliberation Improve an LLM Judge?}
\author{Working research story}
\date{\today}

\begin{document}
\maketitle

\begin{abstract}
Large language models are increasingly used as judges, and reasoning models
can spend hundreds or thousands of tokens before emitting a verdict. The
central question in this paper is not whether chain-of-thought is generally
good or bad. It is whether visible deliberation makes a judge's stated
criterion causally control its decision. Across factual controls, HelpSteer2,
and SummEval, our current Qwen3-8B results support a criterion-use
dissociation. On factual and exact-verification comparisons, the target,
current choice, and final choice are strongly recoverable from residual
states. On evaluative comparisons, the active criterion is often represented
and the model can produce criterion-fluent rationales, but the
criterion-implied target is weakly represented and ordinary rationale text
does not reliably move the verdict. Decisions become reliably
criterion-sensitive when the criterion is operationalized into concrete
comparative evidence or an explicit target. Activation patching and rationale
replay both point to the same bottleneck: awareness of the criterion is not
the same as constructing and using the decision variable implied by that
criterion.
\end{abstract}

\section{One-Sentence Claim}

LLM judges can know what criterion they are supposed to apply, and can
generate rationales that sound criterion-aware, while their final decisions
remain controlled by a different or under-operationalized decision variable.
Operationalized evidence and explicit targets repair this gap far more
reliably than longer or more structured chain-of-thought.

\section{Why This Is the Paper}

Chain-of-thought has its clearest reported benefits on mathematical and
symbolic tasks \citep{sprague2024tocot}, and can be neutral or harmful in
other settings \citep{liu2024mindstep}. LLM judges are a useful stress test
because they often evaluate pairs under plural criteria: helpfulness,
correctness, coherence, relevance, fluency, and overall quality. Existing
LLM-as-judge work already shows sensitivity to order, prompt framing, and
model identity \citep{zheng2023judging,wang2023position,
panickssery2024selfrecognition}. Our contribution is to ask where the
decision actually comes from when a reasoning judge is told to deliberate
under a specific criterion.

The paper should not be framed as ``CoT is unfaithful'' in the broad sense.
The narrower and stronger framing is:

\begin{quote}
  Evaluative CoT can expose a representation/action gap. The model represents
  the requested criterion, but unless the criterion is made operational as
  target evidence, that representation need not control the verdict.
\end{quote}

This framing connects CoT faithfulness concerns \citep{lanham2023faithfulness}
to LLM-as-judge transparency. It also explains why factual verification looks
different: factual tasks provide a goal state that is easier to operationalize.

\section{Result Inventory}

\subsection{Factual-to-Evaluative Motivation}

The broad budget sweep motivates the main contrast. At a 2048-token thinking
budget, factual and exact-verification tasks show high validity and strong
decision convergence. ARC Challenge reaches natural-validity $1.00$,
forced-target success $0.967$, and order-consistent majority rate $1.00$.
GSM8K verification reaches forced-target success $1.00$ and order-consistent
majority rate $1.00$. TruthfulQA reaches forced-target success $0.911$ and
order-consistent majority rate $0.967$.

The factual activation baseline establishes the scale of a well-formed
decision variable. Pair-held-out readouts on factual tasks decode active
criterion at approximately $0.96$--$1.00$ balanced accuracy, criterion target
at $0.82$--$0.93$, current choice at $0.98$--$0.99$, and final choice at
$0.87$--$0.99$. This is the comparison point for the weaker evaluative
readouts.

\subsection{HelpSteer2 Criterion Confirmation}

The locked HelpSteer2 confirmation suite separates criterion language from
operationalized evidence. Criterion-only conditions are modest: early
criterion forced-target adoption is $0.594$ and late criterion adoption is
$0.531$. Supplying criterion-specific score evidence raises the corresponding
conditions to approximately $0.990$ for both early and late evidence. A late
explicit target reaches $1.000$.

The pair-bootstrap effects are large. Evidence operationalization improves
forced-target adoption over criterion-only prompting by $+0.396$ early and
by $+0.458$ late. It improves order-consistent target adoption by $+0.521$ early
and $+0.604$ late. The current behavioral readout therefore supports
``awareness without reliable operationalization'' more strongly than
``late commitment blocks later correction.'' The model can follow the target
when the target is supplied or made computable; it often fails to derive that
target from the criterion alone.

\subsection{HelpSteer2 Structure Controls}

Prompted and computationally enforced structure do not rescue criterion use.
In the structured-CoT suite, free CoT reaches forced-target adoption $0.635$
and order-consistent target adoption $0.521$. The criterion scaffold reaches
$0.594$ and $0.521$, generic scaffold reaches $0.573$ and $0.375$, and
score-evidence reaches $1.000$ and $1.000$.

The adherence analysis shows that the criterion scaffold does change the
surface reasoning. Its criterion-procedure score exceeds free CoT by $+0.120$,
and content-compliance increases by $+0.135$. That does not translate into
better criterion adherence. In fact, both-orders content-compliant traces show
lower order-consistent target adoption in this small sample.

The enforced-structure suite strengthens the negative result. Long free CoT
reaches forced-target adoption $0.625$ and order-consistent adoption $0.500$.
Prompted long reasoning reaches $0.562$ and $0.417$. Enforced generic staging
falls to $0.354$ and $0.292$, and enforced criterion staging reaches only
$0.417$ and $0.208$. Mandatory decomposition increases natural validity but
does not make the criterion govern the decision.

\subsection{HelpSteer2 Mechanistic Evidence}

Fixed-layer readouts show a separation between criterion identity and decision
control. The active criterion is strongly decodable at the intended
criterion-update point, including balanced accuracy $1.000$ at
phase1\_readout\_128 and approximately $0.81$ at phase2\_readout\_0. In
contrast, criterion-target readouts remain weak, roughly $0.39$--$0.50$ across
the main readout points, and final-choice readouts are only modest.

The patching result is the causal centerpiece. A matched criterion-update
residual does not rescue the verdict; its target-selected effect is negative
in the current sample ($-0.111$ versus baseline). Evidence-operationalization
patching is large and selective: matched evidence deltas increase target
selection by $+0.611$, target probability by $+0.604$, and target logit margin
by $+13.289$. Explicit-target patching has the same qualitative pattern,
raising target selection by $+0.625$ and target probability by $+0.601$.

This gives the strongest local mechanistic interpretation: the model contains
states that causally move the decision when the target is operationalized, but
the criterion-update state alone is not sufficient.

\subsection{SummEval External Validation}

SummEval validates the same phenomenon in a different evaluative domain. The
criteria are coherence, consistency, fluency, and relevance. Overall endpoint
adoption is high because many pairs are same-target or tie-to-choice cases:
free CoT reaches forced-target adoption $0.823$ and order-consistent adoption
$0.771$; generic scaffold reaches $0.844$ and $0.802$; criterion scaffold
reaches $0.807$ and $0.719$.

The important stratum is criterion\_flip, where the criterion changes which
summary should win. In that stratum, free CoT reaches only $0.547$
forced-target adoption and $0.469$ order-consistent adoption. Criterion
scaffold is similar or worse at $0.531$ and $0.375$. Generic scaffold is
$0.594$ and $0.500$. Score evidence reaches $1.000$ and $1.000$, while
explicit target reaches $0.938$ and $0.938$.

The endpoint effects make the same point as HelpSteer2. Criterion scaffold
does not improve over free CoT; its forced-target effect is $-0.016$ overall
and its order-consistent effect is $-0.052$. Score evidence improves over free
CoT by $+0.177$ forced adoption and $+0.229$ order-consistent adoption.

\subsection{SummEval Activation Evidence}

The SummEval activation readouts reproduce the representation/use split.
Active criterion is near-perfectly decodable at phase1\_readout\_128 and
phase2\_readout\_0 across conditions. Even by phase2\_readout\_384, active
criterion remains decodable, for example $0.917$ in criterion scaffold and
$0.854$ in free CoT.

Criterion-target readouts are much weaker: roughly $0.54$--$0.59$ at
phase2\_readout\_0 and $0.53$--$0.61$ at phase2\_readout\_384. Current-choice
and final-choice readouts are noisy and much less reliable than the factual
baseline. Thus the model represents which criterion is active, but the
criterion-implied target is not comparably installed as a decision variable.

\subsection{SummEval Rationale Replay}

Rationale replay asks whether text that looks like reasoning can transplant a
decision. Ordinary rationale text does not. Same-free and same-scaffold replay
produce little or no improvement over baseline. Opposite-free and
opposite-scaffold replay also produce little or no movement toward the donor
decision.

Score-evidence replay does move the decision. Same-score replay improves
updated-target adoption by $+0.312$ in native replay and $+0.188$ in quoted
replay. Opposite-score replay reduces recipient target adoption by $-0.531$ in
native replay and $-0.500$ in quoted replay, with probability shifts of
$-0.466$ and $-0.415$. Compared with opposite ordinary rationales, opposite
score evidence shifts recipient-minus-donor logit margin by roughly
$-4.2$ to $-4.4$.

This is the clearest behavioral test of decision origin. The model is not
strongly following arbitrary transplanted rationale prose. It follows
operational comparative evidence.

\section{Paper Structure}

\paragraph{1. Introduction.}
Introduce LLM judges as a deployment-relevant setting where reasoning traces
are used for transparency. State the central risk: a judge may verbalize the
criterion without making the criterion control the verdict.

\paragraph{2. Conceptual model.}
Separate four objects: active criterion, criterion-implied target, current or
final choice, and visible rationale. Deliberation requires these to align over
time. Rationalization-like behavior occurs when criterion information is
present but the choice remains controlled elsewhere. Operationalization is
the process that turns a criterion into concrete target evidence.

\paragraph{3. Factual controls.}
Show that on factual and exact-verification tasks, target and choice are
recoverable and deliberation often converges. This prevents the paper from
claiming that probes are generally weak or that the model cannot represent
decision variables.

\paragraph{4. HelpSteer2 criterion confirmation.}
Present criterion-only, evidence, and explicit-target conditions. Emphasize
that score evidence rescues both early and late updates, while criterion
language alone does not.

\paragraph{5. Mechanistic HelpSteer2 evidence.}
Present active-criterion decoding, weak criterion-target decoding, and the
patching contrast. This is the local causal section.

\paragraph{6. Structured reasoning controls.}
Show that prompting the model to reason in a criterion-aware format, and even
forcing multi-call decomposition, does not reproduce the evidence/target
rescue. This rules out the simple remedy ``ask for a better rationale.''

\paragraph{7. SummEval validation.}
Show the same pattern on coherence, consistency, fluency, and relevance. Focus
on criterion-flip cases. Include SummEval activation readouts.

\paragraph{8. Rationale replay.}
Use ordinary-rationale versus score-evidence transplant as the final
adjudication. Rationale prose does not reliably carry the decision; operational
evidence does.

\paragraph{9. Discussion.}
Frame the result as a representation-to-control gap in evaluative reasoning.
Discuss implications for LLM-as-judge use, CoT faithfulness, rubric design,
and mechanistic audits of evaluators.

\section{Figure and Table Plan}

\begin{description}[leftmargin=3.6cm,style=nextline]
\item[Figure 1]
Conceptual pipeline: criterion identity $\rightarrow$ target construction
$\rightarrow$ choice control $\rightarrow$ rationale. Mark the observed gap
between criterion identity and target/choice control.

\item[Figure 2]
Factual-to-evaluative motivation. Plot factual baseline target/choice
recoverability and behavioral convergence against evaluative conditions.

\item[Figure 3]
HelpSteer2 confirmation behavior. Bars for criterion-only, evidence, and
explicit-target conditions, early and late.

\item[Figure 4]
HelpSteer2 activation and patching. One panel for active-criterion versus
criterion-target decoding; one panel for criterion-update, evidence, and
explicit-target patch effects.

\item[Figure 5]
Structured-reasoning controls. Free CoT, criterion scaffold, generic scaffold,
enforced staging, score evidence, and explicit target.

\item[Figure 6]
SummEval criterion-flip validation. Free/scaffold/generic versus score
evidence and explicit target.

\item[Figure 7]
Rationale replay. Ordinary rationale transplant versus score-evidence
transplant, same and opposite donor directions, native and quoted modes.

\item[Table 1]
Dataset and condition inventory, including pair counts, criteria, transition
strata, branches, and readout budgets.

\item[Table 2]
Result ledger with exact numerical endpoints and confidence intervals for the
main claims.

\item[Appendix tables]
Adherence analysis, enforced-structure token use, dataset-stratified SummEval
effects, and all control patch settings.
\end{description}

\section{Thirty-Billion-Scale Replication Gate}

The larger-model run should be a coherent confirmation suite, not another
exploratory grid. The minimum 30B replication should rerun:

\begin{enumerate}[leftmargin=0.7cm]
\item HelpSteer2 confirmation behavior on the locked pair set:
early criterion, late criterion, early evidence, late evidence, and explicit
target.
\item SummEval structured behavior on the locked SummEval pair set:
free CoT, generic scaffold, criterion scaffold, score evidence, and explicit
target, preserving criterion\_flip, same\_target, and tie\_to\_choice strata.
\item SummEval rationale replay:
same/opposite free CoT, same/opposite criterion scaffold, same/opposite score
evidence, in native and quoted modes.
\item Optional if runtime permits: SummEval activation capture at the same
readout points, focused on criterion\_flip cases rather than the full grid.
\end{enumerate}

The replication succeeds if the qualitative ordering is stable:

\begin{itemize}[leftmargin=0.7cm]
\item score evidence and explicit target remain substantially above
criterion-only and scaffold conditions;
\item criterion scaffold does not reliably outperform free CoT or generic
structure on criterion-flip cases;
\item ordinary rationale replay has little effect compared with baseline;
\item score-evidence replay moves decisions strongly in the donor direction;
\item active criterion, if captured, remains easier to decode than
criterion-implied target or final choice in evaluative settings.
\end{itemize}

The 30B run does not need to replicate enforced structure, adherence scoring,
or full patching before a first conference submission. Those are useful
appendix or follow-up results. The main larger-model question is whether the
behavioral and replay dissociation survives scale.

\section{Venue Positioning}

The natural positioning is interpretability and transparency for LLM judges,
with evaluation as the application domain. The paper is strongest when it is
not sold as a new benchmark and not sold as a universal claim about all
reasoning. It is a mechanistically motivated study of when stated evaluative
criteria become decision-controlling variables.

\begin{description}[leftmargin=3.3cm,style=nextline]
\item[COLM]
Best overall fit if the paper emphasizes language-model reasoning,
judge behavior, and evaluation methodology.
\item[ACL/EMNLP Findings]
Plausible if framed around LLM-as-judge reliability, rubric following, and
evaluation transparency.
\item[BlackboxNLP or NLP interpretability workshops]
Very natural fit, especially if the first version is compact and
mechanistic.
\item[ICLR/NeurIPS workshops]
Good fit for the activation patching and CoT faithfulness angle.
\item[FAccT]
Possible but secondary; the current evidence is more mechanistic than
sociotechnical unless the deployment-risk framing is expanded.
\item[ICLR/NeurIPS/ICML main track]
Ambitious. Plausibility improves if the 30B replication is clean and the
paper is framed as a general representation-to-control failure mode in
deliberative evaluators.
\end{description}

\section{Claim Ledger}

\begin{description}[leftmargin=3.6cm,style=nextline]
\item[Established by current 8B results]
Criterion identity can be represented without reliable criterion-target
construction or choice control. Operational score evidence and explicit targets
repair the decision far more reliably than ordinary or scaffolded CoT.

\item[Supported mechanistically]
Evidence and explicit-target residual directions causally move HelpSteer2
verdicts in patching. Criterion-update residuals alone do not show the same
selective effect.

\item[Supported by replay]
Ordinary rationale text is weak as a decision transplant. Score-evidence text
is strong, including in opposite-target replay and quoted replay.

\item[Not yet claimed broadly]
The phenomenon is not yet shown for all models, all evaluator tasks, or all
reasoning-trained systems.

\item[Main missing gate]
One larger-model replication, preferably Qwen3-30B-A3B or another same-family
reasoning model, showing the same behavioral and replay ordering.

\item[Paper-safe headline]
``In evaluative LLM judging, chain-of-thought can make criteria explicit in
the rationale without making them decision-controlling; operationalized
evidence, not rationale structure, is what reliably redirects the verdict.''
\end{description}

\bibliographystyle{plainnat}
\bibliography{references}
\end{document}
